\section{Resultados}
\label{sec:resultados}

Los resultados se presentan en seis subsecciones: el estudio de calibración del
directorio (el estudio principal de esta versión), el protocolo completo, los
parámetros nativos, la formación del directorio, la capacidad de llenado y el
hemisferio visual, más la medición del costo computacional.
Salvo indicación en contrario, los experimentos usan los parámetros de la
Sección~\ref{sec:parametros}; todos los números provienen de corridas posteriores a la
auditoría de fidelidad (Sección~\ref{sec:auditoria}).

% ---------------------------------------------------------------
\subsection{Calibración del directorio a ocho agentes}
\label{sec:calibracion}

El estudio principal evalúa nueve condiciones de lectura del directorio transactivo
sobre el banco de 411 consultas, con $N \in \{50, 100, 200, 400\}$ consultas de
formación y 5 semillas por celda.
La Tabla~\ref{tab:calibracion} muestra los resultados en $N = 400$ (promedio de 5
semillas); la Figura~\ref{fig:mature-vs-n} muestra la curva completa de la fase madura
y la Figura~\ref{fig:heatmap-condicion} el mapa completo condición $\times$ $N$.

\begin{table}[H]
    \centering
    \caption{Precisión bajo distintas condiciones de lectura del directorio
    ($N = 400$, media de 5 semillas, banco de 411 consultas).
    La precisión temprana no depende de la lectura del directorio, por lo que solo
    varía en las condiciones que alteran las consultas (D, G).
    F es inoperante en esta versión: los labels polisémicos que curaba
    (Apple~Inc.) ya no existen en el vocabulario por masa asociativa, así que
    $F \equiv A$.}
    \label{tab:calibracion}
    \small
    \begin{tabular}{llccc}
        \toprule
        Condición & Descripción & Temprana (\%) & Fidelidad (\%) & Madura (\%) \\
        \midrule
        A (cruda)  & Score de activación sin normalizar             & 88.0 & 87.3 & 80.0 \\
        B1         & División entre conteo de registros por agente  & 88.0 & 93.5 & \textbf{93.2} \\
        B2         & División entre raíz cuadrada del conteo        & 88.0 & 95.0 & 89.0 \\
        C          & Cap de número máximo de registros por agente   & 88.0 & 91.5 & 86.5 \\
        D          & Consultas balanceadas entre clases             & 88.6 & 89.6 & 82.8 \\
        E32        & Cuantización en 32 niveles                     & 88.0 & 91.3 & 84.8 \\
        E64        & Cuantización en 64 niveles                     & 88.0 & 92.0 & 86.5 \\
        F          & Vocabulario curado ($\equiv$ A en v4)          & 88.0 & 87.3 & 80.0 \\
        G (D+B1+F) & Combinación                                    & 88.6 & 93.6 & 92.9 \\
        \bottomrule
    \end{tabular}
\end{table}

\paragraph{B1 sigue siendo la corrección irreducible.}
Igual que en el sistema de tres agentes, la lectura cruda del directorio (A) pierde
13 puntos frente a la lectura normalizada por conteo (80.0\% frente a 93.2\%), y
ninguna condición que manipula síntomas (B2, C, D, E32, E64) la alcanza.
La combinación G no supera a B1 sola (92.9\% frente a 93.2\%): una vez corregida la
causa (evidencia por registro en vez de evidencia total), el balanceo de consultas y la
curación de vocabulario no aportan nada adicional.

\begin{figure}[H]
    \centering
    \includegraphics[width=0.72\textwidth]{fig02_mature_accuracy_vs_N.png}
    \caption{Precisión madura en función de $N$ por condición de lectura.
    B1 es la mejor lectura en todo el rango; la cruda (A) queda sistemáticamente por debajo.}
    \label{fig:mature-vs-n}
\end{figure}

\begin{figure}[H]
    \centering
    \includegraphics[width=0.9\textwidth]{fig03_heatmap_condition_N.png}
    \caption{Mapa de calor de precisión madura por condición experimental y número de
    consultas $N$.}
    \label{fig:heatmap-condicion}
\end{figure}

\paragraph{Desglose por dominio.}
Con la lectura B1 en $N = 400$, la precisión madura por dominio fue:
\texttt{cup} y \texttt{tomato} 100\%, \texttt{car} 98\%, \texttt{cow} y \texttt{dog}
96\%, \texttt{apple} 88\%, \texttt{pear} 86\% y \texttt{horse} 82\%
(Figura~\ref{fig:dominio}).
Las clases de fruta comparten el léxico genérico (\textit{fruit}, \textit{red},
\textit{green}, \textit{round}, \textit{tree}), lo que explica que concentren los
errores residuales.

\begin{figure}[H]
    \centering
    \includegraphics[width=0.85\textwidth]{fig04_domain_accuracy_N80.png}
    \caption{Precisión madura por dominio y condición ($N = 400$).}
    \label{fig:dominio}
\end{figure}

\paragraph{El sesgo hacia \texttt{apple} quedó esencialmente resuelto.}
Las corridas tempranas de 8 clases mostraban al agente \texttt{apple} capturando
$\approx$20\% de las victorias maduras (frente al ideal de 12.5\%) y a \texttt{pear}
colapsando por debajo del 55\%.
Con el encoder determinista oficial y el vocabulario por masa asociativa, la
participación de \texttt{apple} bajo la lectura B1 baja a \textbf{13.0\%} y
\texttt{pear} se recupera a 86\% \emph{sin cambiar el vocabulario}
(Figura~\ref{fig:winner}).
La misma figura muestra que bajo la lectura cruda (A) \texttt{apple} todavía captura el
24\% de las victorias: el residuo que persiste es el sesgo de lectura que B1 corrige,
y con B1 ningún agente se desvía más de 2.5 puntos del ideal.
El análisis de este hallazgo se desarrolla en la Sección~\ref{sec:disc-dominancia}.

\begin{figure}[H]
    \centering
    \includegraphics[width=0.85\textwidth]{fig09_winner_distribution_N80.png}
    \caption{Distribución de agentes ganadores en fase madura por condición ($N = 400$).
    La línea discontinua marca el ideal balanceado de $1/8 = 12.5\%$.}
    \label{fig:winner}
\end{figure}

\paragraph{Estado del directorio.}
La entropía del directorio transactivo alcanza 2.955 bits en la condición A y 2.961
bits en G, sobre un máximo teórico de 3.000 bits ($\log_2 8$): la especialización queda
balanceada al 98.5\% incluso sin corrección, porque el desequilibrio de masas es un
sesgo de \emph{lectura}, no de estructura.
La fidelidad temprana-madura con B1 es del 93.5\%.

% ---------------------------------------------------------------
\subsection{Protocolo completo}
\label{sec:proto-completo}

El protocolo completo (fase temprana $\to$ formación del directorio $\to$ fase madura
con el TME apagado) se corrió sobre el subconjunto de 80 consultas (10 por clase).

\paragraph{Fase temprana.}
El sistema alcanzó una \textbf{precisión temprana del 80.0\%} con una tasa de rechazo
del 5.0\%.
El rechazo no proviene de un filtro de vocabulario: cada token con representación
fastText entra como pista, y es la propia memoria la que rechaza cuando ningún agente
obtiene soporte.

\paragraph{Fase madura.}
Con la lectura B1 la \textbf{precisión madura fue del 82.5\%}; con lectura cruda,
68.8\%.
La fidelidad temprana-madura fue del \textbf{90.0\%}.
Los conteos del directorio al final del experimento fueron
$[50, 16, 44, 25, 30, 14, 17, 31]$ para
(\texttt{apple}, \texttt{car}, \texttt{cow}, \texttt{cup}, \texttt{dog}, \texttt{horse},
\texttt{pear}, \texttt{tomato}), con una entropía de \textbf{2.866 bits} sobre un máximo
de 3.000.
Nótese que la fase madura sobre este subconjunto (82.5\%) es inferior a la del banco
completo (93.2\%): con solo 10 consultas por clase, el directorio dispone de menos
evidencia por dominio; la diferencia mide el efecto del volumen de interacciones, no un
cambio de mecanismo.

\paragraph{Réplica con las 16 consultas del pipeline oficial.}
La réplica con las 16 \texttt{TEST\_QUERIES} (2 por dominio) enruta correctamente las
ocho clases, con conteos $[12, 6, 6, 7, 5, 5, 4, 7]$; los errores residuales
corresponden a consultas de fruta con léxico compartido
(``sweet fruit from an orchard tree'' $\to$ \texttt{apple}).

\paragraph{Sonda de rechazo fuera de dominio.}
Como diagnóstico separado del benchmark, se presentaron 12 consultas
\emph{representables} (todas producen vectores fastText reales) que describen objetos
fuera de los ocho dominios.
La memoria rechazó \textbf{10 de 12} por containment puro (puntuación cero en los ocho
agentes), sin regla léxica alguna.
Las 2 restantes rutean porque comparten tokens con labels legítimos de algún dominio
(p.ej.\ \textit{food}): la frontera entre ``fuera de dominio'' y ``comparte conceptos
con el dominio'' es semántica, no léxica, y el diagnóstico la mide en lugar de esconderla.

% ---------------------------------------------------------------
\subsection[Parámetros nativos iota x kappa]{Parámetros nativos $\iota \times \kappa$}
\label{sec:param-iota-kappa}

Se realizó un barrido de cuadrícula con $\iota \in \{0, 0.25, 0.5, 1.0\}$ y
$\kappa \in \{0, 0.5, 1.0, 1.5\}$ sobre el subconjunto de 80 consultas.
Los resultados se muestran en la Tabla~\ref{tab:grid-iota-kappa}.

\begin{table}[H]
    \centering
    \caption{Precisión temprana y tasa de rechazo en función de $\iota$ y $\kappa$.
    Las columnas de $\kappa$ son idénticas para cada fila, por lo que se muestra una
    sola columna representativa.}
    \label{tab:grid-iota-kappa}
    \begin{tabular}{cccc}
        \toprule
        $\iota$ & $\kappa$ (cualquier valor) & Precisión (\%) & Rechazo (\%) \\
        \midrule
        0.00  & 0 -- 1.5 & 80.0 &   5.0 \\
        0.25  & 0 -- 1.5 & 82.5 &   5.0 \\
        0.50  & 0 -- 1.5 & 65.0 &  14.0 \\
        1.00  & 0 -- 1.5 &  0.0 & 100.0 \\
        \bottomrule
    \end{tabular}
\end{table}

\paragraph{Efecto de $\iota$.}
Con la cuantización por magnitud, una poda ligera ($\iota = 0.25$) mejora marginalmente
la fase temprana (82.5\% frente a 80.0\%) al eliminar algo de ruido de baja
frecuencia, pero la mejora no se traslada de forma sustancial a la fase madura
(B1: 83.8\% frente a 82.5\%) y el patrón destructivo reaparece de inmediato:
$\iota = 0.5$ pierde 15 puntos con un rechazo del 14\%, y $\iota = 1.0$ rechaza todo.
La poda elimina conocimiento escaso pero válido antes de eliminar relaciones espurias.

\paragraph{Efecto de $\kappa$.}
El parámetro de umbral $\kappa$ no cambió ningún resultado en el rango probado: las
proyecciones superan la media del tensor en órdenes de magnitud, de modo que el umbral
relativo nunca es vinculante.

\paragraph{Ninguna condición nativa sustituye a B1.}
En todas las celdas del grid, la lectura madura cruda queda por debajo de la lectura B1
(68.8\% frente a 82.5\% en la línea base; 77.5\% frente a 83.8\% en la mejor
condición nativa).
La calibración alternativa por división entre la media de la memoria
($\div$\texttt{mem.mean}) \emph{empeora} la fase temprana (76.2\%) y distorsiona los
conteos del directorio: con el llenado por instancias las masas de contenido ya están
igualadas por construcción y esa corrección sobra.

% ---------------------------------------------------------------
\subsection{Formación del directorio transactivo}
\label{sec:formacion}

Para estudiar cómo el orden de las interacciones afecta la formación del directorio, se
evaluó la precisión madura completa en un \textit{snapshot} tras cada interacción $k$,
usando cuatro condiciones de presentación (intercalado, barajado con 5 semillas,
bloqueado por dominio, y control con puntuación cruda) sobre el subconjunto de 80
consultas.

\paragraph{Hallazgo principal: a ocho agentes, 10 consultas por clase no bastan.}
A diferencia del sistema de tres agentes (donde el orden intercalado cruzaba el umbral
del 90\% en $k \approx 13$ interacciones), \textbf{ninguna condición alcanza el
umbral del 90\%} con este volumen de interacciones: todas convergen al 82.5\% final.
No es un fallo de convergencia sino de evidencia disponible: con 8 dominios y consultas
que comparten léxico, 10 interacciones por dominio no acumulan suficiente evidencia
distintiva (el mismo protocolo con el banco completo alcanza 93.2\%,
Sección~\ref{sec:calibracion}).

\paragraph{El orden afecta la velocidad, no el destino.}
Todas las condiciones convergen exactamente al mismo estado final: conteos
$[50, 16, 44, 25, 30, 14, 17, 31]$ y entropía 2.866 bits, independientemente del orden
de presentación.
Esta invariancia del destino replica el resultado de tres agentes y confirma que el
directorio acumula evidencia conmutativa: el orden solo determina cuándo se vuelve
utilizable, no qué se aprende.

% ---------------------------------------------------------------
\subsection{Capacidad de llenado}
\label{sec:capacidad}

Se varió el número de imágenes registradas por clase $N \in \{25, 50, 100, 200, 328\}$
utilizando llenado imagen-por-imagen (\textit{imagen-etiqueta}).
La aceptación se determinó por containment inverso con $\xi = 0$.
Los \textit{falsos} son imágenes de otras clases presentadas al agente.

\begin{table}[H]
    \centering
    \caption{Métricas de capacidad en función del número de imágenes registradas $N$
    por clase.
    ``Enrutamiento'' es la precisión del directorio de labels; ``Variedad $L_1$'' es la
    distancia media $L_1$ entre recuperaciones distintas; $H$ es la entropía de la
    memoria de contenido.}
    \label{tab:capacidad}
    \small
    \begin{tabular}{ccccccc}
        \toprule
        $N$ & Train propio (\%) & Test propio (\%) & Falsos (\%) & Variedad $L_1$ & Enrutamiento (\%) & $H$ (bits) \\
        \midrule
         25  & 100 &  0 & 0 & 1.91 & 61.3 & 3.648 \\
         50  & 100 &  0 & 0 & 1.87 & 71.2 & 4.369 \\
        100  & 100 &  2 & 0 & 1.74 & 78.8 & 4.950 \\
        200  & 100 & 20 & 0 & 1.94 & 81.2 & 5.204 \\
        328  & 100 & 56 & 0 & 2.18 & 81.2 & 5.318 \\
        \bottomrule
    \end{tabular}
\end{table}

\paragraph{Generalización por cobertura.}
Con $\xi = 0$ el sistema no generaliza desde cero; la generalización a imágenes de test
aparece a medida que el llenado cubre más del espacio latente de la clase (del 0\% en
$N = 25$ al 56\% en $N = 328$), sin necesidad de ajustar $\xi > 0$ en el contenido.

\paragraph{Especificidad perfecta.}
La tasa de falsos aceptados es del \textbf{0\% en todo el rango} de $N$ probado, ahora
con siete clases contrarias por agente en lugar de dos: el containment inverso mantiene
la discriminación exacta también a ocho clases.

\paragraph{Enrutamiento y entropía.}
El enrutamiento de labels crece con la cobertura y se estabiliza en torno al 81\%
desde $N = 200$; la entropía de la memoria de contenido crece de forma logarítmica con
$N$, coherente con la acumulación de variedad real.
La \textbf{zona de operación} se mantiene en $N = 200$, que además reserva
\texttt{train[200:328]} para las interacciones visuales.

% ---------------------------------------------------------------
\subsection{Hemisferio visual}
\label{sec:hemisferio-visual}

El hemisferio visual opera en dos fases: formación (registro de percepciones reales de
\texttt{train[200:328]} en los directorios visuales) y evaluación (enrutamiento de las
656 imágenes de test mediante el directorio visual por agente).

\paragraph{Fase A: formación con 1024 imágenes.}
Se presentaron 1024 imágenes (128 por clase) en orden intercalado.
Sobre las imágenes aceptadas, la precisión del agente visual fue del \textbf{99.9\%},
con un rechazo del \textbf{0.8\%} (4.2\% con el llenado de 4 variantes; $\approx$50\%
en la versión de tres agentes) y una percepción rechazada antes de la memoria por el
umbral de energía. Cada percepción aceptada registra sus 16 variantes: el registro
completo del TME quedó en $[2048, 2016, 2016, 2016, 2016, 2032, 2048, 2048]$ con
entropía de \textbf{3.000 bits}: el máximo teórico, especialización visual
perfectamente balanceada.

\paragraph{Fase B: enrutamiento de imágenes de test.}
Las 656 imágenes de test se enrutaron mediante los directorios visuales por agente
($M^R_{\mathrm{dir}}$) con lectura B1 estricta ($\xi = 0$, directorios perspectivales):
cada imagen entra por un agente no especialista, que la redirige o la rechaza.
El resultado fue \textbf{631/656 = 96.2\%} de imágenes correctamente asignadas, con un
rechazo del 3.5\% (23 imágenes, cuatro de ellas antes de la memoria por el umbral de
energía) y \textbf{dos falsos enrutamientos} (0.3\%): sobre las imágenes aceptadas, la
precisión es del \textbf{99.7\%}. Los dos falsos son dos vistas del mismo objeto
(\texttt{horse7}, frontal y 035-135) enviadas a \texttt{dog}, que las contiene.
Con el llenado de 4 variantes de la versión anterior el ruteo era del 73.6\% (75.0\%
con $\xi = 2$), con un rechazo del 26.4\% y cero falsos: el directorio no se
equivocaba porque rechazaba una de cada cuatro imágenes.

\paragraph{Justificación de $\xi = 2$.}
El barrido $\xi \in \{0, \ldots, 3\}$ mostró que $\xi = 2$ es el punto dulce de
cobertura, con saturación en $\xi = 3$ y especificidad intacta en todo el rango
(0 falsos sobre sondas sintéticas fuera de dominio).
El análisis del residual es informativo: 625 de las 656 imágenes de test tienen
\emph{cero} coordenadas sin soporte marginal, de modo que el rechazo restante no viene
de huecos por coordenada (que $\xi$ tolera) sino del \emph{containment conjunto}: la
proyección AND no encuentra celdas que sobrevivan a las 64 coordenadas simultáneamente.
Por eso $\xi$ satura, y la palanca para ese residual es densificar el llenado, no
aumentar la tolerancia. El experimento 13 lo midió y esta versión lo adopta: con 16
variantes por imagen en el contenido y por percepción en el directorio, el ruteo de
test sube del 73.6\,\% al 96.2\,\% y la evocación del 85.3 al 97.4\,\%, con dos falsos
de 656 (Sección~\ref{sec:limitaciones}).

\paragraph{Evocación de labels: imagen $\to$ palabras.}
Sobre las imágenes de evaluación aceptadas, la evocación inversa (recuperación $\beta$
desde el latente, decodificada por similitud coseno top-3) produjo al menos una etiqueta
del vocabulario de la clase real en \textbf{113/116 = 97.4\%} de los casos (85.3\% con
el llenado de 4 variantes).

\paragraph{El rechazo y la salida semántica los decide la memoria.}
Reconocimiento, rechazo y reconstrucción son operaciones de la propia memoria: el
encoder ResNet18 solo convierte la imagen en una pista latente, y a partir de ahí la
EAM decide.
La Figura~\ref{fig:rechazo-imagenes} ilustra el rechazo por containment: imágenes
reales que la memoria contiene se aceptan y enrutan, mientras que imágenes reales no
contenidas y las sondas sintéticas de ruido y de píxeles barajados reciben puntuación
cero en los ocho agentes y se rechazan, sin que intervenga ningún filtro externo.
La tercera sonda sintética, un lienzo de color sólido uniforme, \emph{sí} obtiene
soporte y se acepta: es la excepción honesta del panel. El experimento 14 localizó la
causa: no es el fondo liso (los sólidos negro, blanco, de colores y del color de fondo
de ETH-80 se rechazan) sino que el encoder colapsa toda entrada sin estructura, el gris
medio, un objeto desenfocado o sin contraste, en un mismo latente cerca del origen, y
ese punto cae dentro de la envolvente por coordenada de \texttt{horse}, a tres veces la
distancia intraclase de cualquier instancia registrada.
Esta fuga con entradas degeneradas se retoma en la Sección~\ref{sec:limitaciones}.
La Figura~\ref{fig:imagen-etiquetas} muestra la salida semántica: para cada imagen
aceptada, la reconstrucción es lo que la memoria \emph{evoca} (decodificación de la
recuperación de $M^R_{\mathrm{dom}}$), no la imagen de entrada, y el sistema devuelve
las etiquetas asociadas.
Que la reconstrucción difiera de la entrada es esperado y honesto: refleja que decide
la distribución acumulada en la memoria.

\begin{figure}[H]
    \centering
    \includegraphics[width=0.85\textwidth]{fig15_image_rejection.png}
    \caption{Rechazo de imágenes por la memoria (containment, sin filtro léxico).
    Arriba: imágenes reales aceptadas y enrutadas. En medio: imágenes reales no
    contenidas, rechazadas ($\xi = 0$). Abajo: sondas sintéticas fuera de dominio:
    el ruido y los píxeles barajados se rechazan; el color sólido uniforme se acepta,
    una fuga con entradas degeneradas que se discute en la
    Sección~\ref{sec:limitaciones}. El rechazo emerge de la ausencia de soporte, y la
    figura muestra también su excepción.}
    \label{fig:rechazo-imagenes}
\end{figure}

\begin{figure}[p]
    \centering
    \includegraphics[height=0.78\textheight,keepaspectratio]{fig16_image_to_labels.png}
    \caption{Salida semántica imagen $\to$ etiquetas (dos ejemplos por clase). Por cada
    imagen aceptada: entrada (pista), reconstrucción evocada por la memoria (no la
    entrada) y las etiquetas recuperadas, con marca de coincidencia con el vocabulario
    de la clase real.
    Se muestran tanto aciertos como fallos de evocación, coherentes con el 97.4\%
    agregado.}
    \label{fig:imagen-etiquetas}
\end{figure}

% ---------------------------------------------------------------
\subsection{Línea base: una sola EHAM con las ocho clases}

Para medir qué aporta la partición en especialistas, se construyó una
EHAM única (hetero y dos homo) con los mismos 6400 pares etiqueta--latente
que reciben los ocho especialistas y se evaluó con las mismas operaciones
(\texttt{recognize\_gated}, \texttt{recall\_from\_left}), sin ruteo. Se
compararon tres brazos: la EHAM única, los especialistas con el agente
elegido por la verdad de terreno (oráculo, un techo) y el protocolo
transactivo completo, sobre las 171 consultas reservadas, diez semillas y
cortes de 25 a 200 imágenes por clase. Los especialistas reconstruidos son
bit a bit los modelos oficiales.

Con 200 imágenes por clase, la EHAM única acierta la clase el 97.4\,\%
(juez: vecino real más cercano) frente al 97.5\,\% del protocolo y al
100\,\% del oráculo; pero su respuesta cae a 26.7 de la instancia real más
cercana frente a 22.8 y 22.2, y la brecha crece con el contenido (2.0 con
100 registros por clase, 4.5 con 800). Cuando la primera pista reconocida es
una etiqueta compartida entre clases (\emph{fruit}, \emph{animal},
\emph{red}), la EHAM única acierta la clase el 44\,\% y su patrón mezcla
clases (73\,\% de las coordenadas de una sola clase); los especialistas,
100\,\%. En imagen\,$\to$\,texto la EHAM única responde al 96\,\% de las
imágenes de prueba frente al 79\,\% del especialista, y responde con un
dominio ajeno el 9\,\% frente al 1\,\%; la cobertura de ambas crece con el
contenido (el especialista responde el 0, 5, 45 y 79\,\% con 25, 50, 100 y
200 imágenes por clase; la EHAM única, el 5, 54, 81 y 96\,\%). Fuera de
dominio, sobre 40 consultas, el protocolo acepta 4 y la EHAM única 3; con la
doble compuerta de contenido (el agente destino debe contener la pista) el
protocolo acepta 2, como el oráculo. El protocolo queda 2.5 puntos por
debajo del oráculo en todos los cortes: el precio del directorio. Tablas
completas en \texttt{results/experimento11/}.

\subsection{Costo computacional}
\label{sec:costo}

Se midió el tiempo de pared (\textit{wall-clock}) de reproducir el experimento completo
de ocho clases desde cero en una laptop modesta: Intel Core i5 de 13\textsuperscript{a}
generación (12 hilos lógicos, PyTorch usa 8), 16~GB de RAM, \textbf{sin uso de GPU}
(PyTorch en modo CPU; la tarjeta NVIDIA presente no participa).
La Tabla~\ref{tab:costo} resume los tiempos.

\begin{table}[H]
    \centering
    \caption{Costo computacional de la reproducción completa a ocho clases (CPU, sin
    GPU). El entrenamiento del encoder concentra el costo; las decisiones del sistema de
    memorias son operaciones matriciales sobre vectores cuantizados.}
    \label{tab:costo}
    \small
    \begin{tabular}{lc}
        \toprule
        Componente & Tiempo \\
        \midrule
        Encoder ResNet18 (50 épocas, 2624 imágenes, 8 clases) & 79.9 min \\
        Llenado de memorias (instancias + augmentación, 8 agentes) & 15.7 min \\
        Fase temprana + fase madura & 2.9 min \\
        Hemisferio visual (1024 interacciones + test completo) & 34.0 min \\
        \midrule
        Pipeline completo (etapas 1--8) & \textbf{133.5 min} \\
        Batería de análisis completa (calibración, barridos, figuras) & 293.4 min \\
        \bottomrule
    \end{tabular}
\end{table}

El entrenamiento del encoder visual (la única parte con gradientes y retropropagación)
sigue concentrando la mayor parte del costo del pipeline.
El crecimiento respecto a la versión de tres agentes (26.6 min $\to$ 133.5 min) escala
con los datos ($8/3$ de imágenes en el encoder, $\approx$2.7$\times$) y con la
augmentación del llenado, no con el número de agentes en sí: las operaciones de memoria
siguen siendo álgebra matricial sin entrenamiento y se ejecutan en minutos.
La batería de análisis (miles de corridas de evaluación) es trabajo de medición, no
parte de un despliegue del sistema.

La misma medición, repetida en una segunda laptop (Intel Core Ultra 7 265HX, 20 hilos,
16~GB, NVIDIA RTX PRO 1000 de 8~GB), separa lo que aporta cada componente. Con la GPU
oculta, el pipeline baja de 133.5 a 50.6~min y la batería de 159.9 a 70.5~min: una
ganancia de 2.3--2.6$\times$ que viene solo de la CPU, porque las memorias asociativas
corren en \texttt{numpy} y el registro de una hetero es secuencial. Con la GPU activa
solo cambia el entrenamiento del encoder, de 22.3 a 5.1~min (15.8$\times$ respecto de
la primera laptop), y el pipeline queda en 31.9~min; el llenado, los directorios y el
hemisferio visual no la usan. En esa máquina la parte cara del sistema deja de ser el
entrenamiento y pasa a ser la memoria: llenado y hemisferio visual suman 25 de los
32~min. Tabla completa en \texttt{results/computational\_cost/comparacion\_laptops.md}.
